% =====================================================================
\section{Decentralized Registry}
\label{sec:registry}
% =====================================================================

Content addresses give every declaration a stable,
location-independent identity.  The registry organizes the mutable
metadata---who proved what, what is open, what depends on what---around
this immutable core.

Every formalization in Lean already lives in a Git repository.  The
decentralized registry exploits this directly: each participating
project hosts a \texttt{registry/} folder in its own repo, and the
global registry is the union of all such folders.  No additional
infrastructure is needed---no IPFS, no blockchain, no central server.

\sam{Compare P2P DHT (Sam's approach) vs Git-based registry (ours). Sam found DHT doesn't scale, moved to local-first + DHT for publishing locations only.}

This simplicity is possible because \textbf{mathematical consensus is
trivial}.  Unlike distributed systems that must agree on ordering or
resolve write conflicts, a mathematical registry has a monotonic truth
condition: a proposition is proved if \emph{any} repository contains a
valid proof.  Proofs cannot be retracted---once proved, always proved.
Two repositories proving the same proposition under different names is
not a conflict; it is redundancy, and the content address shows they
state the same fact.  There is nothing to vote on, no quorum to reach.
This observation connects to the broader vision of a pluralistic formal
library where multiple communities contribute to a shared body of
verified knowledge~\cite{kohlhase2017qed}.

In concrete terms:

\begin{itemize}[nosep]
  \item \textbf{Publication} is \texttt{git push}---a project generates
    its \texttt{registry/} folder, commits it, and pushes.
  \item \textbf{Verification} is \texttt{lake build}---Lean's kernel
    type-checks the proofs; anyone can clone and verify independently.
  \item \textbf{Consensus} is trivial: one valid proof is enough.
  \item \textbf{Discovery} is a curated list of known repositories
    (\texttt{sources.json}).
\end{itemize}

\paragraph{Annotation attributes.}
Authors annotate their code with two attributes provided by the
\texttt{CA.Registry} module.
\texttt{@[open\_point]} marks a \texttt{def X : Prop} as an open
problem---a mathematical statement published without a proof.
Applying it to a theorem or a non-\texttt{Prop} definition is a
compile-time error.
\texttt{@[publish]} marks any declaration (theorem, definition, axiom)
for publication to the registry.
Both attributes accept an optional description string and support
retroactive application from a separate file, allowing annotation of
existing code or upstream dependencies without modifying the original
source:

\begin{lstlisting}[caption={Retroactive annotation of existing declarations}]
import CA
import MyProject.Theorems

attribute [open_point "My conjecture"] MyProject.SomeConjecture
attribute [publish "Key lemma"]        MyProject.some_lemma
\end{lstlisting}

\paragraph{Declaration status.}
Status is computed automatically from the Lean environment:

\begin{table}[h]
\centering
\begin{tabular}{@{}lll@{}}
  \toprule
  \textbf{Pattern} & \textbf{Status} & \textbf{Rule} \\
  \midrule
  \texttt{theorem foo : P := proof} & proved & Type-checks, no \texttt{sorry} \\
  \texttt{def P : Prop} with \texttt{@[open\_point]} & open & Explicitly marked \\
  \texttt{theorem bar (h : OpenP) : Q} & conditional & Proved modulo open hypotheses \\
  \bottomrule
\end{tabular}
\caption{Automatic status classification.}
\label{tab:status}
\end{table}

For conditional declarations, the registry records which open addresses
are required---the declarations whose proof would promote the entry
from conditional to proved.

\paragraph{Registry folder.}
The \texttt{registry/} folder is generated locally inside each project.
The \texttt{\#ca\_registry "registry/"} command (or the \texttt{ca
registry} CLI) loads the compiled environment, collects annotated
declarations, computes content addresses, classifies status, and writes
two files:

\begin{itemize}[nosep]
  \item \texttt{declarations.json}---the append-only manifest.  Each
    entry contains the content address, name, module, kind, status,
    pretty-printed type, and type dependencies.
  \item \texttt{meta.json}---project-level metadata (hash level,
    open/proved/conditional counts).
\end{itemize}

Entries are never removed.  Status updates (open $\to$ proved) are
reflected by updating the status field in place; Git history provides
the audit trail.

\paragraph{Status cascade.}
When a proof is submitted for an open point, conditional declarations
downstream may become unconditionally proved:

\begin{enumerate}[nosep]
  \item A previously open address acquires a proof---status becomes
    \texttt{proved}.
  \item All \texttt{conditional} entries whose open dependencies
    included the newly proved address are re-evaluated.
  \item If all of an entry's open dependencies are now proved, it is
    promoted to \texttt{proved}.
  \item This cascades transitively.
\end{enumerate}

Cross-project cascades emerge when a consumer merges multiple
registries: an open point in project~A may be proved in project~B,
promoting conditional entries in project~C, with no coordination
between A, B, and~C.  The dependency graph also makes \emph{impact}
computable: the impact of an open point is the number of conditional
declarations transitively downstream of it.

\paragraph{Sources and discovery.}
Because the registry is fully decentralized, a discovery mechanism is
needed so that resolution tools know where to look.  A curated
\texttt{sources.json} file in the CA repository enumerates known Git
repositories that host a \texttt{registry/} folder:

\begin{lstlisting}[style=json,caption={\texttt{sources.json} format}]
{ "registries": [
    { "url": "https://github.com/leanprover-community/mathlib4",
      "registry_path": "registry/" },
    { "url": "https://github.com/user/my-project",
      "registry_path": "registry/" }
  ]
}
\end{lstlisting}

Adding a project to the ecosystem requires five steps: (1)~add
\texttt{require ca} to the project's \texttt{lakefile.lean};
(2)~annotate declarations with \texttt{@[publish]} or
\texttt{@[open\_point]}; (3)~add \texttt{\#ca\_registry "registry/"}
to a Lean file that imports the annotated modules;
(4)~\texttt{lake build}, commit \texttt{registry/}, and push;
(5)~open a PR adding one entry to \texttt{sources.json}.
Step~5 is the only coordination point; everything else is local to the
project.  The static list is the right starting point---it can later be
replaced by automatic discovery (e.g., a GitHub topic tag or a DHT)
without changing the rest of the architecture.

\paragraph{Current maturity.}  The architecture described here is
implemented but early-stage.  Discovery relies on a manually curated
\texttt{sources.json}, there is no mechanism for revoking or correcting
erroneous entries, and metadata conflicts (e.g., two projects assigning
different status to the same address) are resolved by the consumer at
merge time.  The theoretical elegance of monotonic consensus is real,
but the tooling around it is minimal.


% =====================================================================
